\documentclass[11pt]{article}
\usepackage[utf8]{inputenc}
\usepackage{amsmath, amssymb, amsthm}
\usepackage{mathpartir}
\usepackage{stmaryrd}
\usepackage{xcolor}
\usepackage{geometry}
\usepackage{hyperref}
\usepackage{booktabs}

\geometry{margin=1in}

\title{Vyper Bidirectional Typing Rules}
\author{Based on the Pfenning Recipe from Dunfield \& Krishnaswami}
\date{}

\newcommand{\kw}[1]{\texttt{#1}}
\newcommand{\typerule}[1]{\textsc{#1}}
\newcommand{\subtype}[2]{#1 <: #2}
\newcommand{\validelem}[1]{\mathsf{ValidElem}(#1)}
\newcommand{\validdarray}[1]{\mathsf{ValidDArray}(#1)}
\newcommand{\validkey}[1]{\mathsf{ValidKey}(#1)}
\newcommand{\convertible}[2]{\mathsf{Convertible}(#1, #2)}
\newcommand{\checks}{\Leftarrow}
\newcommand{\synth}{\Rightarrow}

\begin{document}

\maketitle

\begin{abstract}
Bidirectional typing rules for Vyper, derived from the type assignment rules
by applying the Pfenning recipe. This formulation separates typing into two
modes: \emph{checking} ($\checks$), where the type is known and verified, and
\emph{synthesis} ($\synth$), where the type is computed from the term structure.
Following the recipe: introduction forms check, elimination forms synthesize,
and a subsumption rule bridges the two modes.
\end{abstract}

\tableofcontents
\newpage

%==============================================================================
\section{Bidirectional Judgment Forms}
%==============================================================================

\subsection{Notation}

We use two principal judgment forms for expressions:

\begin{center}
\begin{tabular}{ll}
$\Gamma; m \vdash e \checks \tau$ & Expression $e$ \emph{checks against} type $\tau$ \\
$\Gamma; m \vdash e \synth \tau$ & Expression $e$ \emph{synthesizes} type $\tau$ \\
\end{tabular}
\end{center}

For statements (which always have type \kw{None} and may extend the context):

\begin{center}
\begin{tabular}{ll}
$\Gamma; m; \tau_r \vdash s \checks \kw{None}$ & Statement $s$ checks against \kw{None} \\
$\Gamma; m; \tau_r \vdash s \checks \kw{None} \dashv \Gamma'$ & Statement $s$ extends context to $\Gamma'$ \\
\end{tabular}
\end{center}

\subsection{Mode Interpretation}

\begin{itemize}
\item \textbf{Checking ($\checks$)}: The type $\tau$ is an \emph{input}. We verify that $e$ has this type.
\item \textbf{Synthesis ($\synth$)}: The type $\tau$ is an \emph{output}. We compute the type from $e$.
\end{itemize}

\subsection{The Pfenning Recipe}

The bidirectionalization follows these principles:

\begin{enumerate}
\item \textbf{Introduction forms check}: Literals, constructors, and $\lambda$-like forms
      have checking conclusions.
\item \textbf{Elimination forms synthesize}: Operations, subscripts, and function applications
      have synthesizing principal judgments.
\item \textbf{Variables synthesize}: The type is looked up from the context.
\item \textbf{Subsumption bridges modes}: When checking but only synthesis is available,
      we synthesize and verify subtyping.
\end{enumerate}

%==============================================================================
\section{Structural Rules}
%==============================================================================

\fbox{$\Gamma; m \vdash e \synth \tau$} \quad (Synthesis)
\hspace{1cm}
\fbox{$\Gamma; m \vdash e \checks \tau$} \quad (Checking)

\subsection{Variable Rule}

Variables synthesize their type from the context:

\begin{mathpar}
\inferrule*[right=Var$\synth$]
  {(x : \tau) \in \Gamma}
  {\Gamma; m \vdash x \synth \tau}
\end{mathpar}

\subsection{Subsumption Rule}

To check an expression against a type, synthesize and verify subtyping:

\begin{mathpar}
\inferrule*[right=Sub$\checks$]
  {\Gamma; m \vdash e \synth \tau_1 \\ \subtype{\tau_1}{\tau_2}}
  {\Gamma; m \vdash e \checks \tau_2}
\end{mathpar}

\paragraph{Note.} This rule is applied when no direct checking rule matches.
It enables expressions that can only synthesize (variables, eliminations)
to be used in checking positions.

%==============================================================================
\section{Literals (Introduction Forms)}
%==============================================================================

Following the recipe, all introduction forms \textbf{check}. Literals are
canonical introductions for their respective types.

\subsection{Boolean Literals}

\begin{mathpar}
\inferrule*[right=True$\checks$]
  { }
  {\Gamma; m \vdash \kw{True} \checks \kw{bool}}

\inferrule*[right=False$\checks$]
  { }
  {\Gamma; m \vdash \kw{False} \checks \kw{bool}}
\end{mathpar}

\subsection{Numeric Literals}

\begin{mathpar}
\inferrule*[right=IntLit$\checks$]
  {-2^{N-1} \leq n \leq 2^{N-1} - 1}
  {\Gamma; m \vdash n \checks \kw{int}N}

\inferrule*[right=UintLit$\checks$]
  {0 \leq n \leq 2^N - 1}
  {\Gamma; m \vdash n \checks \kw{uint}N}

\inferrule*[right=DecLit$\checks$]
  {-2^{167}/10^{10} \leq d \leq (2^{167}-1)/10^{10}}
  {\Gamma; m \vdash d \checks \kw{decimal}}
\end{mathpar}

\subsection{Address and Bytes Literals}

\begin{mathpar}
\inferrule*[right=Addr$\checks$]
  {\text{checksummed}(a)}
  {\Gamma; m \vdash \kw{0x}a \checks \kw{address}}

\inferrule*[right=BytesM$\checks$]
  {|\mathit{hex}| = M \\ 1 \leq M \leq 32}
  {\Gamma; m \vdash \kw{0x}\mathit{hex} \checks \kw{bytes}M}

\inferrule*[right=BytesLit$\checks$]
  {|s| \leq n}
  {\Gamma; m \vdash \kw{b"}s\kw{"} \checks \kw{Bytes}[n]}

\inferrule*[right=StrLit$\checks$]
  {|s| \leq n}
  {\Gamma; m \vdash \kw{"}s\kw{"} \checks \kw{String}[n]}
\end{mathpar}

%==============================================================================
\section{State and Environment Access}
%==============================================================================

State and environment variables \textbf{synthesize} because their types
are determined by the context or are fixed by the language.

\subsection{State Variables}

\begin{mathpar}
\inferrule*[right=StateRead$\synth$]
  {(x : \tau) \in \Gamma \\ \ell(x) \in \{\kw{storage}, \kw{transient}\} \\ m \neq \kw{pure}}
  {\Gamma; m \vdash x \synth \tau}
\end{mathpar}

\subsection{Message Context}

\begin{mathpar}
\inferrule*[right=MsgValue$\synth$]
  {m = \kw{payable}}
  {\Gamma; m \vdash \kw{msg.value} \synth \kw{uint256}}
\end{mathpar}

\subsection{Block and Transaction Environment}

All environment variables synthesize their fixed types:

\begin{center}
\begin{tabular}{@{}llll@{}}
\toprule
\textbf{Variable} & \textbf{Synthesizes} & \textbf{Variable} & \textbf{Synthesizes} \\
\midrule
\kw{block.coinbase} & \kw{address} & \kw{msg.sender} & \kw{address} \\
\kw{block.number} & \kw{uint256} & \kw{msg.gas} / \kw{msg.mana} & \kw{uint256} \\
\kw{block.timestamp} & \kw{uint256} & \kw{msg.data} & \kw{Bytes} \\
\kw{block.difficulty} & \kw{uint256} & \kw{tx.origin} & \kw{address} \\
\kw{block.prevrandao} & \kw{bytes32} & \kw{tx.gasprice} & \kw{uint256} \\
\kw{block.gaslimit} & \kw{uint256} & \kw{chain.id} & \kw{uint256} \\
\kw{block.basefee} & \kw{uint256} & & \\
\kw{block.blobbasefee} & \kw{uint256} & & \\
\kw{block.prevhash} & \kw{bytes32} & & \\
\bottomrule
\end{tabular}
\end{center}

%==============================================================================
\section{Operations (Elimination Forms)}
%==============================================================================

Following the recipe, operations \textbf{synthesize} their result type.
Operands are \textbf{checked} against the required type when known.

\subsection{Boolean Operations}

\begin{mathpar}
\inferrule*[right=And$\synth$]
  {\Gamma; m \vdash e_1 \checks \kw{bool} \\ \Gamma; m \vdash e_2 \checks \kw{bool}}
  {\Gamma; m \vdash e_1\ \kw{and}\ e_2 \synth \kw{bool}}

\inferrule*[right=Or$\synth$]
  {\Gamma; m \vdash e_1 \checks \kw{bool} \\ \Gamma; m \vdash e_2 \checks \kw{bool}}
  {\Gamma; m \vdash e_1\ \kw{or}\ e_2 \synth \kw{bool}}

\inferrule*[right=Not$\synth$]
  {\Gamma; m \vdash e \checks \kw{bool}}
  {\Gamma; m \vdash \kw{not}\ e \synth \kw{bool}}
\end{mathpar}

\subsection{Arithmetic Operations}

For binary arithmetic, we synthesize the first operand to determine the type,
then check the second operand against it:

\begin{mathpar}
\inferrule*[right=BinArith$\synth$]
  {\Gamma; m \vdash e_1 \synth \tau \\ \Gamma; m \vdash e_2 \checks \tau \\
   \tau \in \mathit{NumericT} \\ \oplus \in \{+, -, *, /\!/, \%\}}
  {\Gamma; m \vdash e_1 \oplus e_2 \synth \tau}

\inferrule*[right=DecDiv$\synth$]
  {\Gamma; m \vdash e_1 \checks \kw{decimal} \\ \Gamma; m \vdash e_2 \checks \kw{decimal}}
  {\Gamma; m \vdash e_1\ /\ e_2 \synth \kw{decimal}}

\inferrule*[right=Neg$\synth$]
  {\Gamma; m \vdash e \synth \tau \\ \tau \in \mathit{SignedT}}
  {\Gamma; m \vdash -e \synth \tau}
\end{mathpar}

\subsection{Exponentiation}

\begin{mathpar}
\inferrule*[right=Pow$\synth$]
  {\Gamma; m \vdash e_1 \synth \tau \\ \Gamma; m \vdash e_2 \checks \tau \\
   \tau \in \mathit{IntegerT} \\ \mu(e_1) = \kw{constant} \lor \mu(e_2) = \kw{constant}}
  {\Gamma; m \vdash e_1\ \kw{**}\ e_2 \synth \tau}
\end{mathpar}

\subsection{Bitwise Operations}

\begin{mathpar}
\inferrule*[right=BitBin$\synth$]
  {\Gamma; m \vdash e_1 \synth \tau \\ \Gamma; m \vdash e_2 \checks \tau \\
   \tau \in \mathit{IntegerT} \cup \{\kw{bytes}M, \kw{flag}\ F\} \\ \odot \in \{\&, |, \wedge\}}
  {\Gamma; m \vdash e_1\ \odot\ e_2 \synth \tau}

\inferrule*[right=BitNot$\synth$]
  {\Gamma; m \vdash e \synth \tau \\ \tau \in \{\kw{uint256}, \kw{bytes32}, \kw{flag}\ F\}}
  {\Gamma; m \vdash \mathord{\sim} e \synth \tau}
\end{mathpar}

\subsection{Shift Operations}

\begin{mathpar}
\inferrule*[right=Shift$\synth$]
  {\Gamma; m \vdash e_1 \synth \tau \\ \Gamma; m \vdash e_2 \synth \tau' \\ \tau' \in \mathit{IntegerT} \\
   \tau \in \{\kw{uint256}, \kw{int256}, \kw{bytes32}\} \\ \diamond \in \{\ll, \gg\}}
  {\Gamma; m \vdash e_1\ \diamond\ e_2 \synth \tau}
\end{mathpar}

\subsection{Comparisons}

Comparisons synthesize \kw{bool}:

\begin{mathpar}
\inferrule*[right=NumCmp$\synth$]
  {\Gamma; m \vdash e_1 \synth \tau \\ \Gamma; m \vdash e_2 \checks \tau \\
   \tau \in \mathit{NumericT} \\ \bowtie\ \in \{<, \leq, >, \geq\}}
  {\Gamma; m \vdash e_1\ \bowtie\ e_2 \synth \kw{bool}}

\inferrule*[right=Eq$\synth$]
  {\Gamma; m \vdash e_1 \synth \tau \\ \Gamma; m \vdash e_2 \checks \tau}
  {\Gamma; m \vdash e_1 == e_2 \synth \kw{bool}}

\inferrule*[right=NEq$\synth$]
  {\Gamma; m \vdash e_1 \synth \tau \\ \Gamma; m \vdash e_2 \checks \tau}
  {\Gamma; m \vdash e_1\ !=\ e_2 \synth \kw{bool}}
\end{mathpar}

\subsection{Membership Tests}

\begin{mathpar}
\inferrule*[right=InSArray$\synth$]
  {\Gamma; m \vdash a \synth \tau[n] \\ \Gamma; m \vdash e \checks \tau}
  {\Gamma; m \vdash e\ \kw{in}\ a \synth \kw{bool}}

\inferrule*[right=InDArray$\synth$]
  {\Gamma; m \vdash a \synth \kw{DynArray}[\tau, n] \\ \Gamma; m \vdash e \checks \tau}
  {\Gamma; m \vdash e\ \kw{in}\ a \synth \kw{bool}}

\inferrule*[right=InFlag$\synth$]
  {\Gamma; m \vdash f \synth \kw{flag}\ F \\ \Gamma; m \vdash e \checks \kw{flag}\ F}
  {\Gamma; m \vdash e\ \kw{in}\ f \synth \kw{bool}}
\end{mathpar}

%==============================================================================
\section{Subscript and Attribute Access}
%==============================================================================

These are elimination forms and thus \textbf{synthesize}.

\subsection{Subscript Access}

\begin{mathpar}
\inferrule*[right=SArrayIdx$\synth$]
  {\Gamma; m \vdash e \synth \tau[n] \\ \Gamma; m \vdash i \synth \tau_i \\ \tau_i \in \mathit{IntegerT} \\
   \mu(i) = \kw{constant} \Rightarrow 0 \leq i < n}
  {\Gamma; m \vdash e[i] \synth \tau}

\inferrule*[right=DArrayIdx$\synth$]
  {\Gamma; m \vdash e \synth \kw{DynArray}[\tau, n] \\ \Gamma; m \vdash i \synth \tau_i \\ \tau_i \in \mathit{IntegerT}}
  {\Gamma; m \vdash e[i] \synth \tau}

\inferrule*[right=MapIdx$\synth$]
  {\Gamma; m \vdash h \synth \kw{HashMap}[\kappa, \tau] \\ \Gamma; m \vdash k \checks \kappa}
  {\Gamma; m \vdash h[k] \synth \tau}

\inferrule*[right=TupleIdx$\synth$]
  {\Gamma; m \vdash t \synth (\tau_0, \ldots, \tau_{n-1}) \\ 0 \leq i < n \\ \mu(i) = \kw{constant}}
  {\Gamma; m \vdash t[i] \synth \tau_i}
\end{mathpar}

\subsection{Attribute Access}

\begin{mathpar}
\inferrule*[right=StructAttr$\synth$]
  {\Gamma; m \vdash e \synth \kw{struct}\ S \\ (f : \tau) \in S}
  {\Gamma; m \vdash e.f \synth \tau}

\inferrule*[right=AddrAttr$\synth$]
  {\Gamma; m \vdash e \synth \kw{address} \\ m \neq \kw{pure} \\
   (a, \tau_a) \in \{(\kw{balance}, \kw{uint256}), (\kw{codesize}, \kw{uint256}),
           (\kw{codehash}, \kw{bytes32}), (\kw{is\_contract}, \kw{bool})\}}
  {\Gamma; m \vdash e.a \synth \tau_a}
\end{mathpar}

%==============================================================================
\section{Composite Literals (Introduction Forms)}
%==============================================================================

Composite literals are introduction forms and thus \textbf{check}.

\begin{mathpar}
\inferrule*[right=ArrayLit$\checks$]
  {\Gamma; m \vdash e_i \checks \tau \text{ for all } i \in 1..n}
  {\Gamma; m \vdash [e_1, \ldots, e_n] \checks \tau[n]}

\inferrule*[right=Tuple$\checks$]
  {\Gamma; m \vdash e_i \checks \tau_i \text{ for all } i \in 1..n}
  {\Gamma; m \vdash (e_1, \ldots, e_n) \checks (\tau_1, \ldots, \tau_n)}

\inferrule*[right=StructCtor$\checks$]
  {\kw{struct}\ S = \{f_1 : \tau_1, \ldots, f_n : \tau_n\} \\
   \Gamma; m \vdash e_i \checks \tau_i \text{ for all } i}
  {\Gamma; m \vdash S(f_1 = e_1, \ldots, f_n = e_n) \checks \kw{struct}\ S}
\end{mathpar}

%==============================================================================
\section{Conditional Expression}
%==============================================================================

The conditional expression has both checking and synthesis variants.

\subsection{Checking Variant}

When a goal type is known, both branches check against it:

\begin{mathpar}
\inferrule*[right=IfExp$\checks$]
  {\Gamma; m \vdash e_c \checks \kw{bool} \\ \Gamma; m \vdash e_1 \checks \tau \\ \Gamma; m \vdash e_2 \checks \tau}
  {\Gamma; m \vdash e_1\ \kw{if}\ e_c\ \kw{else}\ e_2 \checks \tau}
\end{mathpar}

\subsection{Synthesis Variant}

When no goal type is known, branches must synthesize the same type:

\begin{mathpar}
\inferrule*[right=IfExp$\synth$]
  {\Gamma; m \vdash e_c \checks \kw{bool} \\ \Gamma; m \vdash e_1 \synth \tau \\ \Gamma; m \vdash e_2 \checks \tau}
  {\Gamma; m \vdash e_1\ \kw{if}\ e_c\ \kw{else}\ e_2 \synth \tau}
\end{mathpar}

%==============================================================================
\section{Type Conversion and Zero Initialization}
%==============================================================================

These operations have explicit type annotations and thus \textbf{synthesize}.

\begin{mathpar}
\inferrule*[right=Convert$\synth$]
  {\Gamma; m \vdash e \synth \tau_1 \\ \convertible{\tau_1}{\tau_2}}
  {\Gamma; m \vdash \kw{convert}(e, \tau_2) \synth \tau_2}

\inferrule*[right=Empty$\synth$]
  { }
  {\Gamma; m \vdash \kw{empty}(\tau) \synth \tau}
\end{mathpar}

%==============================================================================
\section{Function Calls}
%==============================================================================

Function calls are elimination forms and \textbf{synthesize} the return type.
Arguments are \textbf{checked} against the expected parameter types.

\subsection{Internal Function Calls}

\begin{mathpar}
\inferrule*[right=InternalCall$\synth$]
  {(f : (\overline{\tau}) \xrightarrow{\kw{internal}, m_f} \tau_f) \in \Gamma \\
   \Gamma; m \vdash e_i \checks \tau_i \text{ for all } i \\
   m_f \leq m}
  {\Gamma; m \vdash f(\overline{e}) \synth \tau_f}
\end{mathpar}

\subsection{External Function Calls}

\begin{mathpar}
\inferrule*[right=ExtCall$\synth$]
  {\Gamma; m \vdash e \synth \kw{interface}\ I \\
   (f : (\overline{\tau}) \xrightarrow{\kw{external}, m_f} \tau_f) \in I \\
   m_f \geq \kw{nonpayable} \\
   \Gamma; m \vdash e_i \checks \tau_i \text{ for all } i \\
   m \geq \kw{nonpayable}}
  {\Gamma; m \vdash \kw{extcall}\ e.f(\overline{e}) \synth \tau_f}

\inferrule*[right=StaticCall$\synth$]
  {\Gamma; m \vdash e \synth \kw{interface}\ I \\
   (f : (\overline{\tau}) \xrightarrow{\kw{external}, m_f} \tau_f) \in I \\
   m_f \in \{\kw{view}, \kw{pure}\} \\ m_f \leq m \\
   \Gamma; m \vdash e_i \checks \tau_i \text{ for all } i}
  {\Gamma; m \vdash \kw{staticcall}\ e.f(\overline{e}) \synth \tau_f}
\end{mathpar}

%==============================================================================
\section{Dynamic Array Methods}
%==============================================================================

\begin{mathpar}
\inferrule*[right=Pop$\synth$]
  {\Gamma; m \vdash a \synth \kw{DynArray}[\tau, n] \\ \mu(a) = \kw{modifiable}}
  {\Gamma; m \vdash a.\kw{pop}() \synth \tau}
\end{mathpar}

%==============================================================================
\section{Statements}
\label{sec:statements}
%==============================================================================

All statements \textbf{check against \kw{None}}. Some statements extend the
typing context.

\fbox{$\Gamma; m; \tau_r \vdash s \checks \kw{None}$}
\hspace{0.5cm}
\fbox{$\Gamma; m; \tau_r \vdash s \checks \kw{None} \dashv \Gamma'$}

\subsection{Method Calls as Statements}

\begin{mathpar}
\inferrule*[right=Append$\checks$]
  {\Gamma; m \vdash a \synth \kw{DynArray}[\tau, n] \\ \Gamma; m \vdash e \checks \tau \\
   \mu(a) = \kw{modifiable}}
  {\Gamma; m \vdash a.\kw{append}(e) \checks \kw{None}}
\end{mathpar}

\subsection{Variable Declaration}

Variable declarations extend the context:

\begin{mathpar}
\inferrule*[right=VarDecl$\checks$]
  {\Gamma; m \vdash e \checks \tau}
  {\Gamma; m \vdash (x : \tau = e) \checks \kw{None} \dashv \Gamma, x : \tau}
\end{mathpar}

\subsection{Assignment}

\begin{mathpar}
\inferrule*[right=Assign$\checks$]
  {(x : \tau) \in \Gamma \\ \Gamma; m \vdash e \checks \tau \\
   \mu(x) = \kw{modifiable} \\
   \ell(x) \in \{\kw{storage}, \kw{transient}\} \Rightarrow m \geq \kw{nonpayable}}
  {\Gamma; m \vdash x = e \checks \kw{None}}

\inferrule*[right=AugAssign$\checks$]
  {(x : \tau) \in \Gamma \\ \Gamma; m \vdash e \checks \tau \\ \tau \in \mathit{NumericT} \\
   \oplus \in \{+, -, *, /\!/, \%, \kw{**}\} \\
   \mu(x) = \kw{modifiable} \\
   \ell(x) \in \{\kw{storage}, \kw{transient}\} \Rightarrow m \geq \kw{nonpayable}}
  {\Gamma; m \vdash x\ \oplus\!\!= \ e \checks \kw{None}}

\inferrule*[right=AugAssignBit$\checks$]
  {(x : \tau) \in \Gamma \\ \Gamma; m \vdash e \checks \tau \\
   \tau \in \mathit{IntegerT} \cup \{\kw{bytes}M, \kw{flag}\ F\} \\
   \odot \in \{\&, |, \wedge\} \\
   \mu(x) = \kw{modifiable} \\
   \ell(x) \in \{\kw{storage}, \kw{transient}\} \Rightarrow m \geq \kw{nonpayable}}
  {\Gamma; m \vdash x\ \odot\!\!= \ e \checks \kw{None}}

\inferrule*[right=AugAssignShift$\checks$]
  {(x : \tau) \in \Gamma \\ \Gamma; m \vdash e \synth \tau' \\ \tau' \in \mathit{IntegerT} \\
   \tau \in \{\kw{uint256}, \kw{int256}, \kw{bytes32}\} \\ \diamond \in \{\ll, \gg\} \\
   \mu(x) = \kw{modifiable} \\
   \ell(x) \in \{\kw{storage}, \kw{transient}\} \Rightarrow m \geq \kw{nonpayable}}
  {\Gamma; m \vdash x\ \diamond\!\!= \ e \checks \kw{None}}
\end{mathpar}

\subsection{Control Flow}

\begin{mathpar}
\inferrule*[right=If$\checks$]
  {\Gamma; m \vdash e \checks \kw{bool} \\
   \Gamma; m; \tau_r \vdash s_1 \checks \kw{None} \\
   \Gamma; m; \tau_r \vdash s_2 \checks \kw{None}}
  {\Gamma; m; \tau_r \vdash \kw{if}\ e\kw{:}\ s_1\ \kw{else:}\ s_2 \checks \kw{None}}

\inferrule*[right=ForRange$\checks$]
  {\Gamma; m \vdash e_1 \synth \tau \\ \Gamma; m \vdash e_2 \checks \tau \\ \tau \in \mathit{IntegerT} \\
   \Gamma, i : \tau; m; \tau_r \vdash s \checks \kw{None}}
  {\Gamma; m; \tau_r \vdash \kw{for}\ i : \tau\ \kw{in range}(e_1, e_2)\kw{:}\ s \checks \kw{None}}

\inferrule*[right=ForSArray$\checks$]
  {\Gamma; m \vdash a \synth \tau[n] \\
   \Gamma, x : \tau; m; \tau_r \vdash s \checks \kw{None}}
  {\Gamma; m; \tau_r \vdash \kw{for}\ x : \tau\ \kw{in}\ a\kw{:}\ s \checks \kw{None}}

\inferrule*[right=ForDArray$\checks$]
  {\Gamma; m \vdash a \synth \kw{DynArray}[\tau, n] \\
   \Gamma, x : \tau; m; \tau_r \vdash s \checks \kw{None}}
  {\Gamma; m; \tau_r \vdash \kw{for}\ x : \tau\ \kw{in}\ a\kw{:}\ s \checks \kw{None}}
\end{mathpar}

\subsection{Return}

\begin{mathpar}
\inferrule*[right=Return$\checks$]
  {\Gamma; m \vdash e \checks \tau_r}
  {\Gamma; m; \tau_r \vdash \kw{return}\ e \checks \kw{None}}

\inferrule*[right=ReturnVoid$\checks$]
  {\tau_r = \kw{None}}
  {\Gamma; m; \tau_r \vdash \kw{return} \checks \kw{None}}
\end{mathpar}

\subsection{Assert}

\begin{mathpar}
\inferrule*[right=Assert$\checks$]
  {\Gamma; m \vdash e \checks \kw{bool}}
  {\Gamma; m \vdash \kw{assert}\ e \checks \kw{None}}

\inferrule*[right=AssertMsg$\checks$]
  {\Gamma; m \vdash e \checks \kw{bool} \\ \Gamma; m \vdash \mathit{msg} \checks \kw{String}[n] \\ n \leq 1024}
  {\Gamma; m \vdash \kw{assert}\ e,\ \mathit{msg} \checks \kw{None}}

\inferrule*[right=AssertUnreach$\checks$]
  {\Gamma; m \vdash e \checks \kw{bool}}
  {\Gamma; m \vdash \kw{assert}\ e,\ \kw{UNREACHABLE} \checks \kw{None}}
\end{mathpar}

\subsection{Terminating Statements}

\begin{mathpar}
\inferrule*[right=Raise$\checks$]
  {\Gamma; m \vdash e \checks \kw{String}[n] \\ n \leq 1024}
  {\Gamma; m \vdash \kw{raise}\ e \checks \kw{None}}

\inferrule*[right=RaiseUnreach$\checks$]
  { }
  {\Gamma; m \vdash \kw{raise}\ \kw{UNREACHABLE} \checks \kw{None}}

\inferrule*[right=Selfdestruct$\checks$]
  {\Gamma; m \vdash e \checks \kw{address} \\ m \geq \kw{nonpayable}}
  {\Gamma; m \vdash \kw{selfdestruct}(e) \checks \kw{None}}

\inferrule*[right=Break$\checks$]
  {\text{inside a loop}}
  {\Gamma; m \vdash \kw{break} \checks \kw{None}}

\inferrule*[right=Continue$\checks$]
  {\text{inside a loop}}
  {\Gamma; m \vdash \kw{continue} \checks \kw{None}}
\end{mathpar}

\subsection{Event Logging}

\begin{mathpar}
\inferrule*[right=Log$\checks$]
  {\kw{event}\ E = \{a_1 : \tau_1, \ldots, a_n : \tau_n\} \\
   \Gamma; m \vdash e_i \checks \tau_i \text{ for all } i \\
   m \geq \kw{nonpayable}}
  {\Gamma; m \vdash \kw{log}\ E(a_1 = e_1, \ldots, a_n = e_n) \checks \kw{None}}
\end{mathpar}

\subsection{Ether Transfer and Low-Level Calls}

\begin{mathpar}
\inferrule*[right=Send$\checks$]
  {\Gamma; m \vdash e_1 \checks \kw{address} \\ \Gamma; m \vdash e_2 \checks \kw{uint256} \\
   m \geq \kw{nonpayable}}
  {\Gamma; m \vdash \kw{send}(e_1, e_2) \checks \kw{None}}

\inferrule*[right=RawCallStmt$\checks$]
  {\Gamma; m \vdash e_1 \checks \kw{address} \\ \Gamma; m \vdash e_2 \checks \kw{Bytes}[n] \\
   \text{no \kw{max\_outsize} specified} \\ m \geq \kw{nonpayable}}
  {\Gamma; m \vdash \kw{raw\_call}(e_1, e_2, \ldots) \checks \kw{None}}
\end{mathpar}

For \kw{raw\_call} with output:

\begin{mathpar}
\inferrule*[right=RawCallExpr$\synth$]
  {\Gamma; m \vdash e_1 \checks \kw{address} \\ \Gamma; m \vdash e_2 \checks \kw{Bytes}[n] \\
   \kw{max\_outsize} = k \\ m \geq \kw{nonpayable}}
  {\Gamma; m \vdash \kw{raw\_call}(e_1, e_2, \ldots, \kw{max\_outsize}=k) \synth \kw{Bytes}[k]}
\end{mathpar}

%==============================================================================
\section{Function Definitions}
%==============================================================================

\begin{mathpar}
\inferrule*[right=FunDef$\checks$]
  {\Gamma, x_1 : \tau_1, \ldots, x_n : \tau_n; m_f; \tau_r \vdash \mathit{body} \checks \kw{None}}
  {\Gamma \vdash \kw{@}v\ \kw{@}m_f\ \kw{def}\ f(\overline{x : \tau}) \to \tau_r \checks
   (\overline{\tau}) \xrightarrow{v, m_f} \tau_r}
\end{mathpar}

%==============================================================================
\section{Summary: Bidirectionalization Principle}
%==============================================================================

\begin{center}
\begin{tabular}{@{}lll@{}}
\toprule
\textbf{Form Category} & \textbf{Direction} & \textbf{Examples} \\
\midrule
Variables & Synthesize ($\synth$) & $x$, state variables \\
Literals (Introduction) & Check ($\checks$) & \kw{True}, $42$, \kw{"hello"} \\
Composite Literals & Check ($\checks$) & $[e_1, \ldots]$, $(e_1, e_2)$, $S(\ldots)$ \\
Operations (Elimination) & Synthesize ($\synth$) & $e_1 + e_2$, $\kw{not}\ e$ \\
Subscript/Attribute & Synthesize ($\synth$) & $e[i]$, $e.f$ \\
Function Calls & Synthesize ($\synth$) & $f(e_1, \ldots)$, \kw{extcall} \\
Type Annotations & Synthesize ($\synth$) & \kw{convert}, \kw{empty} \\
Conditional Expr & Both & $e_1\ \kw{if}\ c\ \kw{else}\ e_2$ \\
Statements & Check $\checks$ \kw{None} & assignments, control flow \\
\bottomrule
\end{tabular}
\end{center}

%==============================================================================
\begin{thebibliography}{9}

\bibitem{dunfield2020}
Jana Dunfield and Neel Krishnaswami.
\textit{Bidirectional Typing}.
arXiv:1908.05839, 2020.

\bibitem{dunfield2004}
Jana Dunfield and Frank Pfenning.
Tridirectional Typechecking.
\textit{POPL}, 2004.

\bibitem{pierce2002types}
Benjamin C. Pierce.
\textit{Types and Programming Languages}.
MIT Press, 2002.

\end{thebibliography}

\end{document}
